\head{EGA V 预备笔记：超平面截面与锥投影}
\source{\texttt{EGAVaiello.pdf} 物理第1--3页；本篇与 \texttt{EGAV5.ps} 的开头部分对应}
\noindent\textbf{底本标题署名：}Alexandre Grothendieck 的 EGA V；其“预备笔记”的翻译与编辑：Piotr Blass、Joseph Blass。

\egasection{编写摘要}
这一稿把应收入最终稿的全部结果杂陈在一起，作了详细概述。要达到最终稿的形式，需要对目前这个“第零阶段”彻底重新组织。第一步大概应当是重订计划（毫无疑问，现在的第11、12、14、15节将大幅提前）。我甚至还没有写第16节；原则上，它应当没有困难，也不以任何方式影响前面的各节，因为涉及的只是转换表述。

你会注意到，这里有一条命题10.3，它应当出现在前面的某个段落。

就此，我还想告诉你，我另有若干彼此很不相同、但都涉及双有理映射的结果，很愿意把它们放到某处。我觉得这些还不足以独立构成一节。你对放在哪里有什么建议吗？我打算不久把它们连同这份笔记的第16节寄给你。\ednote{询问 AG：第16节究竟是否写过。}

此外，目前的第20段大概会扩展成两段，其中一段由 Grassmann 概形上属于“初等几何”类型的结果组成。必要时，是否也可以把我跟你提过的关于双有理变换的补充材料放在这里（如果找不到更合适的地方）？

\egasection{拟定目录}
\begin{enumerate}[label=\arabic*)]
\item 记号。
\item 泛超平面截面：局部性质。
\item 泛超平面截面：几何不可约性与连通性。
\item 变动超平面截面：“充分一般”的截面。
\item Seidenberg 类型定理。
\item 任意超平面截面的连通性。
\item 特殊超平面截面与多截面构造的应用。
\item 例外超平面集合的维数。
\item 射影浸入的改变。
\item 超平面截面的铅笔与爆破簇的纤维化。
\item Grassmann 概形。
\item 前述结果向线性截面的推广。
\item 初等态射与 M. Artin 的一个定理。
\item 锥投影。
\item 部分前述结果的公理化。
\item 转换为线性系的语言。
\end{enumerate}
以上是底本列出的拟定目录，其中部分内容仍属写作计划；本书实际收录范围以卷首和来源记录为准。

